\section{Resultados Preliminares}\label{sec:results}

Os testes iniciais alcançaram uma taxa de acerto próxima a cem por cento no reconhecimento das bolhas, distinguindo com grande precisão as noites normais das noites afetadas pelo fenômeno climático.

Ao utilizarmos as ferramentas de explicação visual, constatamos que a máquina focou sua atenção consistentemente na parte superior do céu, exatamente a região onde as bolhas costumam ascender e se manifestar na realidade. Isso comprova que o modelo aprendeu o padrão físico autêntico.

Além disso, simulamos o funcionamento do modelo de forma ininterrupta durante uma noite inteira. O sistema se comportou de maneira impecável ao processar centenas de fotos limpas consecutivas sem gerar nenhum falso positivo, reforçando sua vocação para atuar no monitoramento constante e autônomo.
